\section{从 Union Bound 到 Markov 与 Chebyshev}
\label{sec:05-Probability/07-Transforms-and-Bounds/02}

你接手一个系统：一千个组件，每个组件一天内故障的概率是 \(0.001\)。系统只要有一个组件故障就算失败。你被要求给出系统一天内失败的概率。

如果各组件独立，\(1-0.999^{1000}\approx 0.632\)。如果它们因共享电源而完全同步，概率就只有 \(0.001\)。现实几乎总是夹在两者之间——你既不能默认独立，也不知道相关的具体结构。精确概率算不出来。

这不是方法的缺陷，是信息的边界。你只知道每个组件的边缘故障概率，却不知道联合分布。这种情况下，你需要的不是一个精确数字，而是一个在任何依赖结构下都成立的上界——一个不会因为你不知道相关性就失效的安全围栏。

\subsection{Union Bound：把全局坏事件拆成局部件}

最粗糙也最普适的工具是 Union Bound（Boole 不等式）。对任意事件 \(A_1,\ldots,A_n\)——不要求独立，不要求互斥：

\[
P\Bigl(\bigcup_{i=1}^n A_i\Bigr) \le \sum_{i=1}^n P(A_i).
\]

证明不需要任何技巧，测度的次可加性是概率公理的一部分。等号成立当且仅当这些事件两两不交——实际中很少满足，所以上界通常是松的。

回到那一千个组件：

\[
P(\text{系统失败}) \le 1000 \times 0.001 = 1.
\]

上界等于一比不说还糟。但如果每个组件的故障概率降到 \(10^{-6}\)，上界就是 \(10^{-3}\)——不需要知道组件之间是否独立，这个保证已经足够用。

Union Bound 的力量在于策略而非精度：把一个难以直接处理的全局坏事件（``至少一个组件坏了''），覆盖成许多容易控制的小事件（``这个组件坏了''）之和。随机算法的正确性证明、高维统计中的联合置信、编码理论中的错误概率分析——到处是这个策略的身影。代价是松：当事件有大量重叠时，上界可以远大于真实概率。

\subsection{Markov 不等式：均值压下尾部}

如果说 Union Bound 管的是事件层面，Markov 不等式则管到了随机变量的量级。条件只有两个：非负，知道均值。

若 \(X\ge 0\) 且 \(a>0\)，则

\[
P(X\ge a) \le \frac{\E[X]}{a}.
\]

证明可以短到一行。考虑逐点不等式：

\[
a\cdot \mathbf{1}_{\{X\ge a\}} \le X.
\]

验证逐点成立：当 \(X\ge a\)，左边等于 \(a\)，右边 \(X\ge a\)，成立；当 \(0\le X<a\)，左边为 \(0\)，右边非负，成立。两边取期望（期望保序）：

\[
a\,P(X\ge a) \le \E[X],
\]

即不等式本身。指示函数把``事件的概率''和``随机变量的期望''焊在一起，这是 Markov 不等式唯一用到的手法。

举个例子把它落地。某服务窗口的等待时间平均 \(5\) 分钟。等待时间非负，所以

\[
P(X\ge 20) \le \frac{5}{20} = \frac14.
\]

等超过二十分钟的概率，不超过四分之一。你可能会说：这个界也太松了——如果等待时间是指数分布，\(P(X\ge 20)=e^{-20/5}=e^{-4}\approx 0.018\)，小了十多倍。你说得对。Markov 只用了均值这一个数，分布的整个形状都被它扔掉了，它当然松。但反过来说，只用均值这一个数能做到的最紧上界，就是 Markov 给的。某些两点分布可以打到等号——\(P(X=a)=\E[X]/a\)，\(P(X=0)=1-\E[X]/a\)（要求 \(\E[X]\le a\)）——说明在仅知非负和均值的条件下，不能普遍做得更好。

非负条件不能摘掉。若 \(X\) 可以取很大的负值，均值可能很小甚至为负，却仍然有相当大的正尾概率——Markov 对此完全沉默，因为它从未为负值设计。

\subsection{Chebyshev：方差为偏离均值画了圈}

Markov 只管单侧的界，而且需要非负。能不能管到``偏离均值''——两侧同时抓？如果能构造一个和偏离有关的非负变量，套上 Markov 即可。

方差正是为此而生。对任意有有限均值 \(\mu\) 和有限方差 \(\sigma^2\) 的随机变量，构造非负变量 \((X-\mu)^2\)，套 Markov：

\[
P(|X-\mu|\ge a) = P\bigl((X-\mu)^2 \ge a^2\bigr) \le \frac{\E[(X-\mu)^2]}{a^2} = \frac{\sigma^2}{a^2}.
\]

这就是 Chebyshev 不等式。用标准差做刻度更干净：令 \(a=k\sigma\)，

\[
P(|X-\mu|\ge k\sigma) \le \frac{1}{k^2}.
\]

所以至少 \(1-1/k^2\) 的概率落在均值周围 \(k\) 个标准差之内。这条规则对任意有方差的分布成立，不要求对称、不要求单峰、不要求 Gaussian。

代入数字感受分布的差异。\(k=2\) 时，Chebyshev 只保证至少 \(75\%\) 落在两倍标准差内。Gaussian 给出约 \(95\%\)，差距来自 Chebyshev 对分布形状的一无所知——它必须同时覆盖 Gaussian 的紧致和 Pareto 的弥散。\(k=3\) 时至少 \(89\%\)，Gaussian 是 \(99.7\%\)。Chebyshev 是一个不能更保守的普适上界，代价是上界的强度也同样保守。

\subsection{单侧版本：用上方向信息可以更紧}

De Chebyshev 是双尾的。如果只关心中一侧——比如只问``利润会不会低于某个阈值''而不关心超额利润——可以利用方向信息拿到更紧的界。

Cantelli 不等式给出：

\[
P(X-\mu\ge a) \le \frac{\sigma^2}{\sigma^2+a^2}.
\]

证明的方法透露出一个可迁移的策略。思路仍然是 Markov，但更主动：先平移，再平方，引入一个自由参数，最后优化它。

对任意 \(u>0\)，

\[
\begin{aligned}
P(X-\mu\ge a)
&= P(X-\mu+u\ge a+u) \\
&\le P\bigl((X-\mu+u)^2 \ge (a+u)^2\bigr) \\
&\le \frac{\E[(X-\mu+u)^2]}{(a+u)^2}
   = \frac{\sigma^2 + u^2}{(a+u)^2}.
\end{aligned}
\]

第一步把单侧事件装进一个更大的平方事件，第三步对非负的 \((X-\mu+u)^2\) 使用 Markov。现在 \(u\) 是我们引入的自由参数——选得好可以收紧界。求右端关于 \(u\) 的最小值，取 \(u=\sigma^2/a\)，代回即得 Cantelli 不等式。

这个``引入自由参数再优化''的动作，正是下一节 Chernoff 界的方法核心。Chernoff 不过是在指数尺度上重复同一策略：用 \(e^{tX}\) 代替这里的平方，然后对 \(t\) 优化。

Cantelli 比直接把 Chebyshev 的双侧界掰一半用在单侧更紧——多用了方向信息，少付了保守性的代价。

\subsection{这些界能证明平均会稳定}

Chebyshev 不只是一个不等式，它可以推出一条定律。设 \(X_1,\ldots,X_n\) 独立同分布，均值 \(\mu\)，方差 \(\sigma^2\)。样本均值 \(\bar X_n\) 的方差为 \(\sigma^2/n\)。Chebyshev 直接上场：

\[
P(|\bar X_n-\mu|\ge\varepsilon) \le \frac{\sigma^2}{n\varepsilon^2} \longrightarrow 0,\qquad n\to\infty.
\]

这就是弱大数定律的有限方差版本——不需要特征函数，不需要 Fourier 分析，只需要一条用方差和 Markov 堆出来的初等不等式。它以 \(1/n\) 的速度衰减，不算快。对有界独立变量，本单元末节的 Hoeffding 集中界会把它提升为关于 \(n\) 的指数衰减，但那是用了更多信息（有界性）换来的升级。

\subsection{信息代价一览}

总结本节四条界的层级：

\begin{itemize}
\tightlist
\item Union Bound：只知道事件的边缘概率，不知道它们的联合——给出和的上界；
\item Markov：只知道非负变量的均值——给出单尾界；
\item Chebyshev：知道均值和方差——给出双尾界，对称但不分方向；
\item Cantelli：在 Chebyshev 的基础上加用方向信息——单尾收紧。
\end{itemize}

每多知道一层信息，界就可以收紧一层。信息越少，承诺越粗——这不是方法有什么毛病，而是信息限制的精确反映。把它反过来想同样成立：如果你手里只有一个均值和方差，就不要指望拿到 Gaussian 级别的尾部精度。那不是不等式的错，是你对那个分布知道的实在太少。
